\subsection{The function-field centralizer and the margin of scan 127}

\paragraph{Statement B7 (Proved: the birational centralizer criterion).}\label{res:B7}
Let $K$ be a field finitely generated over $\mathbf Q$, and let $X/K$ be a smooth, proper, geometrically connected
curve of any genus, and $L=K(X)$. Fix an algebraic closure $\overline L$
and the resulting copy of $\overline K\subset\overline L$. Put
\[
 E_L=\operatorname{Gal}(\overline L/L),\qquad
 P_L=\operatorname{Gal}(\overline L/\overline K(X)).
\]
Let $s:G_K\to E_L$ be a continuous section of the restriction
$E_L\twoheadrightarrow G_K$. Then
\[
 C_{P_L}(s(G_K^\circ))\ne1
\]
holds exactly for the sections whose image lies in the decomposition
group $D_v$ of a prolongation $v$ to $\overline L$ of a $K$-rational place
$x$ of $L$. In that case the prolongation is unique and
\[
 C_{P_L}(s(G_K^\circ))=I_v\simeq\widehat{\mathbf Z}(1),\qquad
 C_{E_L}(s(G_K^\circ))=I_v.
\]
The same assertion holds for a germ over $G_{K'}$, for every finite
extension $K'/K$, with a $K'$-rational place of $K'(X)$ and with unchanged
geometric group $P_L$. This proves the margin formula of cote 149,
\S4, manuscript p.~63, scan 127, in relative transcendence degree one.

\begin{proof}
For every nonempty $K$-open $U\subset X$, restriction gives a continuous
surjection
\[
 q_U:E_L\twoheadrightarrow E_U=\pi_1(U),
 \qquad q_U|_{P_L}:P_L\twoheadrightarrow\pi_1(U_{\overline K}).
\]
For $V\subset U$ the comparison map $q_{VU}:E_V\twoheadrightarrow E_U$
satisfies $q_U=q_{VU}q_V$, and all maps commute with their projections to
$G_K$. The function-field argument in B5 gives
$E_L=\varprojlim_U E_U$ and
$P_L=\varprojlim_U\pi_1(U_{\overline K})$.
Deleting at least three geometric points gives a cofinal family of
hyperbolic affine opens, in every genus.

Write $C=C_{P_L}(s(G_K^\circ))$ and choose $1\ne z\in C$.
Some hyperbolic open $U_0$ satisfies $q_{U_0}(z)\ne1$.
For every hyperbolic $V\subset U_0$, its image $z_V$ remains nontrivial,
because it maps to $z_{U_0}$. Apply the centralizer theorem of Part A and its proved finitely generated extension in Statement AE1 to
$s_V=q_Vs$. It gives a unique cusp $c_V$ of the completed universal
pro-cover of $V$, fixed by $s_V(G_K)$, with
\[
 C_{\pi_1(V_{\overline K})}(s_V(G_K^\circ))=I_{c_V}.
\]

We verify compatibility instead of assuming it. The map of normalized
completions associated to $q_{VU_0}$ sends $c_V$ either to a cusp over
$X\setminus U_0$ or to a point over $U_0$. In the latter case its inertia
maps trivially, since the universal cover of $U_0$ is \emph{\'etale}
there. That would kill $z_{U_0}$, which is impossible. Thus the image is
a cusp, and its inertia contains the nonidentity element $z_{U_0}$.
The distinct-inertia assertion of Stix
\cite[Proposition 11]{B_Stix2008}, for hyperbolic curves in characteristic
zero with full profinite geometric group, identifies that image cusp
with $c_{U_0}$. The same argument for any pair $V\subset U$ in this
cofinal family gives compatibility of all $c_U$.

These compatible points determine a valuation prolongation $v$ of
$\overline L$. Here is the field comparison. Every finite extension of
$L$ inside $\overline L$ is unramified on some $U$; the normalization of
$X$ in that extension is a finite level in the corresponding completed
universal tower. The points supplied by $c_U$ restrict compatibly on all
such finite normalizations. Their valuation rings have a directed union
inside $\overline L$, and for each $a\in\overline L^*$ either $a$ or
$a^{-1}$ belongs to that union, as this holds already at a finite level
containing $a$. This is the required valuation ring. Since all $c_U$ are
fixed by the image of the full section, $s(G_K)\subset D_v$. The
underlying geometric point of $X$ is fixed by $G_K$, so its residue field
is $K$: it is a $K$-rational place $x$.

The characteristic-zero local exact sequence is
\[
 1\longrightarrow I_v\simeq\widehat{\mathbf Z}(1)
 \longrightarrow D_v\longrightarrow G_K\longrightarrow1.
\]
For a hyperbolic $U$ deleting $x$, its map to $E_U$ is the map of the
punctured henselian neighborhood at $x$. Stix
\cite[Lemma 10]{B_Stix2008} proves that this map is injective for
$\chi(U)\le0$; its image is the cusp decomposition group, and its
restriction identifies $I_v$ with $I_{c_U}$. The arithmetic action on this
inertia is the cyclotomic action. Thus $I_v\subset C$.

For the reverse inclusion, let $w\in C$. For every $U$ in the cofinal
family, $q_U(w)\in I_{c_U}$. Choose $a\in I_v$ with
$q_{U_0}(a)=q_{U_0}(w)$. The comparison maps
$I_{c_V}\to I_{c_{U_0}}$ for $V\subset U_0$ are isomorphisms, because
both are identified with the same $I_v$ by the preceding local
injections. Therefore $q_V(a)=q_V(w)$ for every such $V$. The inverse-limit
isomorphism gives $w=a$, so $C=I_v$.

Conversely a section in a decomposition group $D_v$ centralizes $I_v$
on $G_K^\circ$, so $C\ne1$ and the equality just proved applies.
Uniqueness follows from the uniqueness of every $c_U$: any other
prolongation fixed by the section would supply the same compatible
system of points. The equality of the centralizers in $E_L$ and $P_L$
is B6. Passing to a finite $K'/K$ replaces $E_L$ by the inverse image of
$G_{K'}$, and does not change $P_L$; repeating the same argument gives
the assertion for germs.
\end{proof}

\paragraph{Statement B8 (Refuted in the written function-field form;
proved for the model quotient).}\label{res:B8}
For a section in $D_v$ the two groups written as trivial in (32)--(33)
of scan 127 are both $I_v$, and are nontrivial. For a hyperbolic model
$U/K$ containing the rational point $x$ underneath $v$, its point section
$s_x=q_Us$ satisfies
\[
 C_{E_U}(s_x(G_K^\circ))
 =C_{\pi_1(U_{\overline K})}(s_x(G_K^\circ))=1.
\]
The comparison responsible for the two formulas is the continuous
surjection $q_U:E_L\twoheadrightarrow E_U$; it kills $I_v$ when $x\in U$.

\begin{proof}
The function-field formula is B7. For the model, the point section and a
cuspidal section are disjoint by Stix
\cite[Theorem 14(2)]{B_Stix2008}, whose hypothesis is injectivity of point
sections on proper curves of genus at least two over all finite base
extensions. Over the present finitely generated field this hypothesis follows from the
no-divisible-points argument of Part A, extended in Statement AE1, and the Jacobian Kummer
comparison. The centralizer theorem of Part A therefore gives the
geometric equality $1$, and B6 gives the arithmetic equality.
The local map at an interior point is the map from a strict henselian
trait into an \emph{\'etale} neighborhood, so its ramification inertia is
in the kernel of $q_U$.

For a completely explicit counterexample take $K=\mathbf Q$,
$L=\mathbf Q(t)$, $x=0$ on $X=\mathbf P^1$, and
$U=\mathbf P^1\setminus\{1,2,\infty\}$. A parameter $t$ at $0$ splits
$D_v\to G_{\mathbf Q}$, by the local construction in Stix
\cite[Proposition 8]{B_Stix2008}; compose that splitting with
$D_v\hookrightarrow E_L$. The inertia $I_v$ is already seen to be
nontrivial on the quadratic extension $\overline{\mathbf Q}(t^{1/2})$
of $\overline{\mathbf Q}(t)$, where it maps onto $\mathbf Z/2\mathbf Z$.
The program \texttt{B\_field\_certificates.py} checks the discriminant
$\operatorname{disc}_T(T^2-t)=4t$, with a simple zero at $t=0$.
All of $I_v$ is killed on passage to $E_U$, and the section there is the
point section of $0\in U(\mathbf Q)$.
\end{proof}

We inspected scan 127, including its lower half. The two written
subscripts read $L$. The proved comparison supports the replacement of
the \emph{mathematical assertion} by its $U$ version in the indicated
context. It does not authorize silently changing the transcription, and
does not establish what the author intended. The slash/comma variation
in $\pi_{L/K}$ has no effect on the argument once that group is defined
as the displayed geometric kernel. Notice also that the arithmetic and
geometric formulas with $L$ are still equivalent to each other by B6;
their common value, rather than that equivalence, is the erroneous
assertion of triviality for a decomposition section.

\paragraph{Statement B9 (Proved: the precise logical content of statement 2).}\label{res:B9}
For $L=K(X)$ as in B7, the following two predicates on the set of germs
of sections have the same extension:
\[
 \{[s]:C_{P_L}(s(G_{K'}^\circ))\ne1\}
 =
 \{[s]:s(G_{K'})\subset D_v\text{ for a rational place }v\}.
\]
Here the equivalence relation on sections is conjugation by $P_L$, and
$K'/K$ ranges through finite extensions. Consequently the universal
assertion in scan 45, restricted to relative transcendence degree one,
and the birational section assertion a) of scans 126--127 over all finite
base extensions are logically equivalent.

\begin{proof}
Statement B7 proves equality of the predicates for each actual section
over each $K'$. Both predicates are invariant under $P_L$-conjugacy:
conjugation sends the centralizer to its conjugate and sends a valuation
decomposition group to that of the conjugate prolongation. Restriction
of a section to an open subgroup gives the finite-extension case of the
same assertion. Universal quantification over those sections therefore
gives the stated equivalence. We make no assertion that the universal
predicate is satisfied by all sections.
\end{proof}

The exact remaining assertion is the existence, for every birational
section $s$, of a nonidentity element of
$C_{P_L}(s(G_K^\circ))$. Our proved criterion detects every section for
which such an element is present; it does not supply the element for an
arbitrary section. A basic unproved instance of that remaining assertion
is $K=\mathbf Q$, $L=\mathbf Q(t)$, with an arbitrary continuous section
$G_{\mathbf Q}\to G_{\mathbf Q(t)}$. The finite quotient calculations
in this work concern specified decomposition sections and do not settle
that universal instance.

\subsection{Packets and recovery of the curve's points and open models}

\paragraph{Statement B10 (Proved: packets are the nonzero centralizer fibres).}\label{res:B10}
For the same $K,X,L$, let $\mathcal S_L(K)$ be the set of continuous
sections $G_K\to E_L$, modulo $P_L$-conjugacy. On its subset
$\mathcal S_L^{\rm dec}(K)$ of decomposition sections the map
\[
 [s]\longmapsto[C_{P_L}(s(G_K^\circ))]_{P_L}
\]
has one fibre for each $x\in X(K)$. For a fixed prolongation $v$ above
$x$, that fibre is precisely the injective image of the set of splittings
$G_K\to D_v$, modulo $I_v$-conjugacy. It is a torsor under
\[
 H^1(G_K,I_v)=H^1(K,\widehat{\mathbf Z}(1))
 \simeq\varprojlim_n K^*/(K^*)^n.
\]
Every $x\in X(K)$ has a nonempty such fibre. This proves the packet
description of cote 149, \S3, manuscript p.~28, scan 57, for function
fields of curves over fields finitely generated over $\mathbf Q$.

\begin{proof}
The map is well defined by conjugacy functoriality, and B7 identifies
its nontrivial value with $I_v$. We first prove the required normalizer
identity in the function field:
\[
 N_{E_L}(I_v)=D_v,\qquad N_{P_L}(I_v)=I_v.
\]
For $g$ in the first normalizer, the inertia groups of $v$ and $gv$
coincide. On every sufficiently small hyperbolic open deleting their
centres they have equal nontrivial images. Stix's Proposition 11 then
identifies the corresponding cusps. Varying the open and the finite
normalization identifies the two valuations, by the construction in B7.
Thus $g\in D_v$. Conversely $D_v$ normalizes its inertia by definition.
Intersecting with $P_L$ gives the second identity, since the residue
field of the prolongation in $\overline L$ is $\overline K$.

Now suppose two splittings $s,t$ into $D_v$ become conjugate by
$g\in P_L$ in $E_L$. By B7, $gI_vg^{-1}=I_v$. The normalizer identity
forces $g\in I_v$, giving injectivity. Distinct underlying points have
nonconjugate inertia groups: a $P_L$-conjugation acts over
$X_{\overline K}$, and hence cannot change their underlying geometric
points. Conversely all prolongations above the same $K$-rational point
are conjugate by $P_L$. At each finite geometric Galois normalization the
Galois group acts transitively on the points above that point; compactness
of $P_L$ gives a compatible conjugating element over the full tower.
This proves the description of the fibres.

For completeness, fix one splitting $s_0$ into $D_v$. Every other one is
uniquely $s_a(\gamma)=a(\gamma)s_0(\gamma)$, where $a$ is a continuous
$I_v$-valued $1$-cocycle for the cyclotomic action. The equation expressing
that $s_a$ is a homomorphism is exactly the cocycle equation, and
conjugation by $I_v$ is exactly addition of a continuous coboundary.
Thus the set is a torsor under continuous $H^1(G_K,I_v)$. The Kummer
comparison in Part C(a) gives the displayed inverse limit. A local
parameter at $x$ gives the initial splitting: the compatible covers
obtained by adjoining its $n$th roots give the splitting of the punctured
henselian neighborhood, as in the opened Stix Proposition 8. Its
inclusion into $E_L$ provides the required decomposition section.
\end{proof}

\paragraph{Statement B11 (Proved: what the packet construction recovers).}\label{res:B11}
From $E_L\to G_K$ with its given cyclotomic character, the nontrivial
centralizers of germs of sections recover $X(\overline K)$ as a $G_K$-set
and hence recover $X(K')$ for every finite extension $K'/K$.
They also recover, for every finite $G_K$-stable subset
$I\subset X(\overline K)$, the quotient extension
$E_{X\setminus I}\to G_K$. These data determine the $K$-isomorphism
class of every open model $X\setminus I$, with compatibility under
inclusions. No assertion that every birational section is a decomposition
section is needed for this reconstruction: the nontrivial-centralizer
subset is already exactly characterized by B7.

\begin{proof}
For a finite $K'/K$ inside $\overline K$, form the set of packets
\[
 \mathscr P_L(K')=
 \bigl\{[C_{P_L}(s(G_{K'}^\circ))]_{P_L}:
            s\text{ a section over }G_{K'},\ C_{P_L}(s(G_{K'}^\circ))\ne1\bigr\}.
\]
This is the set of centralizer conjugacy classes, rather than the set
of individual splittings in its torsor fibres. B10 identifies it with
$X(K')$. Restriction to $G_{K''}$ for $K''\supset K'$ leaves the
centralizer $I_v$ unchanged by B7, so it defines the transition map on
these packet sets and corresponds exactly to $X(K')\subset X(K'')$.
Every geometric point is defined over a finite extension, and every
such point has a decomposition section by B10. Hence their filtered
colimit is exactly $X(\overline K)$.

For $\tau\in G_K$, choose $g\in E_L$ mapping to $\tau$. The map
$[I_v]_{P_L}\mapsto[gI_vg^{-1}]_{P_L}$ is independent of that lift,
since two lifts differ by $P_L$. It sends the packet over $K'$ to the
packet over $\tau(K')$ at the conjugate geometric point. This defines
the required $G_K$-action on the colimit. Its fixed points under
$G_{K'}$ are exactly $X(K')$.

Write $U=X\setminus I$. In $E_L$, form the closed normal subgroup
\[
 R_U=\overline{\langle\!\langle I_v:
       v\text{ lies over }x\in U(\overline K)\rangle\!\rangle}.
\]
It lies in $P_L$. We claim that the quotient map identifies
$E_L/R_U$ with $E_U$. To check this it suffices to check its finite
quotients. A finite Galois extension of $L$ kills these inertia groups
exactly when its normalization over the regular curve $U$ is unramified
at every closed point. In characteristic zero the residue extensions are
separable; a finite normalization unramified at these discrete valuation
rings is finite \emph{\'etale}. Conversely a finite \emph{\'etale} cover
has trivial inertia there. Thus the two profinite quotients have exactly
the same finite quotients, with the same maps from $E_L$. Their canonical
continuous comparison is an isomorphism. This also proves compatibility:
for $V\subset U$ one has $R_V\subset R_U$, producing precisely the
continuous surjection $E_V\twoheadrightarrow E_U$.

To recover the schemes, we use the opened Tamagawa
\cite[Theorem (0.3), with Conjecture (0.1)(ii) as its stated formulation]
{B_Tamagawa1997}. Its exact setting is a field $k$ finitely generated over
$\mathbf Q$ and two smooth, geometrically connected, affine hyperbolic
$k$-curves. The natural map from their $k$-isomorphisms to continuous
isomorphisms of their arithmetic fundamental groups over $G_k$, modulo
inner automorphisms from the target geometric group, is bijective.

Choose a finite $G_K$-stable set $J\supseteq I$ containing at least three
geometric points. Then $V=X\setminus J$ is affine hyperbolic, and the
reconstructed quotient $E_V\to G_K$ determines $V$ by that theorem.
Its unique smooth proper compactification determines $X$. The remaining
labels $J\setminus I$ are its cusp inertia classes, already supplied by
the quotient maps from $E_L$, and adding exactly these points to $V$
gives $U$. For two choices $J_1,J_2$, let $J=J_1\cup J_2$ and first
reconstruct the common hyperbolic refinement $V=X\setminus J$.
Fill the cusp points labelled by $J\setminus J_i$ in its smooth proper
compactification. The arithmetic fundamental group of the resulting
open is the quotient of $E_V$ by the normal closure of those cusp
inertias: the same finite-normalization argument used for $R_U$ proves
this assertion. That quotient is the already reconstructed
$E_{X\setminus J_i}$. The result of filling is itself affine
hyperbolic, so Tamagawa's bijection gives the unique isomorphism to
the independently reconstructed $X\setminus J_i$ inducing this
specified group isomorphism. A curve isomorphism extends uniquely to
the smooth proper compactifications, and its maps on cusp inertia
classes give the specified labels there. Adding the further labelled
points needed for $U$ now gives compatible isomorphisms of the desired
opens. On three choices the isomorphisms agree on their common
hyperbolic refinement by the same uniqueness. Thus both the recovered
isomorphism class and the inclusion maps are independent of the
auxiliary set $J$.
\end{proof}
